\documentclass[11pt,a4paper]{report}
\usepackage[utf8]{inputenc}
\usepackage[italian]{babel}
\usepackage[margin=1in]{geometry}
\usepackage{hyperref}
\usepackage{xcolor}

\definecolor{primary}{RGB}{38, 66, 139}

\begin{document}

\chapter*{Esercizi Lezione 6: Real-time & WebRTC}

\section*{Esercizio 1: WebSocket Chat}
Implementa un server Express che utilizza la libreria \texttt{ws} per gestire una chat di gruppo.
\begin{enumerate}
    \item Il server deve mantenere una lista di client connessi.
    \item Ogni messaggio ricevuto da un client deve essere inviato in broadcast a tutti gli altri.
    \item Crea un client in Vanilla JS che permetta di inviare messaggi e visualizzare quelli ricevuti.
\end{enumerate}

\section*{Esercizio 2: WebRTC P2P Loopback}
Basandoti sull'esempio in \texttt{src/lezione6/3.video-streaming/}:
\begin{enumerate}
    \item Aggiungi un pulsante per "Mute" dell'audio locale senza interrompere il video.
    \item Modifica il codice per gestire l'evento di chiusura della connessione (\textit{onconnectionstatechange}) e pulire correttamente la UI.
\end{enumerate}

\section*{Esercizio 3: Analisi di Rete}
\begin{itemize}
    \item Utilizza gli Strumenti per Sviluppatori del browser (scheda \textit{Network}) per osservare l'handshake HTTP di una connessione WebSocket.
    \item Identifica gli header \texttt{Upgrade: websocket} e \texttt{Connection: Upgrade}.
\end{itemize}

\end{document}
